\section{可插拔后端：\texttt{heap\_ops\_t} 接口}

与第~\ref{ch:pmm}~章的 \texttt{pmm\_ops\_t} 相同模式——接口与实现解耦：

\codefile{src/include/kernel/kmalloc.h}
\begin{lstlisting}[style=cstyle]
typedef struct heap_ops {
    const char* name;
    void  (*init)(void* start, size_t size);
    void* (*alloc)(size_t size);
    void  (*free)(void* ptr);
    void  (*stats)(kheap_stats_t* out);
} heap_ops_t;
\end{lstlisting}

\begin{figure}[H]
  \centering
  % figures/ch10/fig02-kmalloc-dispatch.tex — auto-extracted from chapter source
\begin{tikzpicture}[
  box/.style={draw, rounded corners=3pt, minimum height=0.7cm,
              minimum width=3.2cm, font=\small, align=center},
  arr/.style={-{Stealth[length=5pt]}, thick},
]
  \node[box, fill=blue!8]   (caller) at (0, 0)    {调用者 (VMA / Shell)};
  \node[box, fill=yellow!12](api)    at (0, -1.3)  {\texttt{kmalloc()} / \texttt{kfree()}};
  \node[box, fill=orange!12](disp)   at (0, -2.6)  {\texttt{kmalloc.c} dispatch};
  \node[box, fill=green!12] (ff)     at (-2.8,-4.0){\texttt{heap\_first\_fit.c}};
  \node[box, fill=green!12] (slab)   at (2.8,-4.0) {\texttt{heap\_slab.c}};
  \draw[arr] (caller)--(api);
  \draw[arr] (api)--(disp);
  \draw[arr] (disp)--(ff);
  \draw[arr] (disp)--(slab);
\end{tikzpicture}

  \caption{\texttt{kmalloc} 可插拔后端调用链}
  \label{fig:kmalloc-dispatch}
\end{figure}

\texttt{kconfig.h} 编译期选择后端（\texttt{0}=first\_fit, \texttt{1}=slab），
\texttt{kmain.c} 启动时注册并调用 \texttt{kmalloc\_init()}。

\begin{thinkbox}
"dispatch + 后端"模式在 Linux 内核中随处可见（VFS、I/O 调度器、协议栈）。
如果不用函数指针表，而在每个调用点用 \texttt{\#if} 选择实现，会有什么问题？
（提示：代码膨胀与维护成本）
\end{thinkbox}

% ============================================================================

